\subsection{INTERSECTING HYPERPLANES}

Recall (Algebra, Chap. II, §9, no. 3) that two affine subspaces P and P' of E are said to be parallel if there exists a vector t in T such that \(P' = t + P\) . It is clear that the relation ``P and P' are parallel'' is an equivalence relation on the set of affine subspaces of E.

Lemma 1. Any two non-parallel hyperplanes have non-empty intersection.

Let H and \(H'\) be two non-parallel hyperplanes, \(a \in H\) and \(a' \in H'\) ; there exist two hyperplanes M and \(M'\) of the vector space T such that \(H = M + a\) and \(H' = M' + a'\) ; since H and \(H'\) are not parallel, we have \(M \neq M'\) and hence \(T = M + M'\) ; hence there exist \(u \in M\) and \(u' \in M'\) such that \(a' - a = u - u'\) , and the point \(u + a = u' + a'\) belongs to \(H \cap H'\) .

Lemma 2. Let H and H' be two distinct hyperplanes of E, and f, \(f'\) two affine functions on E such that H (resp. H') consists of the points a in E such that \(f(a) = 0\) (resp. \(f'(a) = 0\) ). Finally, let L be a hyperplane of E. Assume that one of the following hypotheses is satisfied:

\begin{enumerate}
\def\labelenumi{\alph{enumi})}
\item
  The hyperplanes H, H' and L are parallel.
\item
  The hyperplanes H and \(\mathrm{H}'\) are not parallel, and \(\mathrm{H} \cap \mathrm{H}' \subset \mathrm{L}\) .
\end{enumerate}

Then there exist real numbers \(\lambda, \lambda'\) , not both zero, such that \(L\) consists of those points \(a \in E\) at which the affine function \(g = \lambda.f + \lambda'.f'\) vanishes.

The lemma being trivial when \(\mathbf{L} = \mathbf{H}\) , we assume that there exists a point \(a\) in \(\mathbf{L}\) with \(a \notin \mathbf{H}\) . Put \(\lambda = f'(a)\) , \(\lambda' = -f(a)\) and

\[
g = \lambda . f + \lambda^ {\prime}. f ^ {\prime};
\]

then \(\lambda' \neq 0\) since \(a \notin \mathrm{H}\) ; moreover, since \(\mathrm{H} \neq \mathrm{H}'\) , there exists \(b \in \mathrm{H}\) such that \(b \notin \mathrm{H}'\) , so that \(f(b) = 0\) , \(f'(b) \neq 0\) , and thus \(g(b) = -f(a).f'(b)\) is nonzero. The set \(\mathbf{L}_1\) of points where the affine function \(g \neq 0\) vanishes is then a hyperplane of \(\mathbf{E}\) ; we have \(g(a) = 0\) , so \(a \in \mathbf{L}_1\) .

\begin{enumerate}
\def\labelenumi{\alph{enumi})}
\item
  Assume that H and \(H'\) are parallel: g and f both vanish at every point of \(L_{1} \cap H\) , hence so does \(f'\) since \(\lambda' \neq 0\) ; thus every point of \(L_{1} \cap H\) belongs to \(H'\) ; but since H and \(H'\) are parallel and distinct, they are disjoint, so \(L_{1} \cap H = \varnothing\) , and Lemma 1 shows that \(L_{1}\) is parallel to H. Since \(a \in L\) and \(a \in L_{1}\) , we thus have \(L = L_{1}\) .
\item
  Assume that H and \(\mathbf{H}'\) are not parallel: by Lemma 1, there is a point \(c\) in \(\mathbf{H} \cap \mathbf{H}'\) ; we give E the vector space structure obtained by taking \(c\) as the origin. Then \(H \cap H'\) is a vector subspace of E of codimension 2, and since \(a \notin H\) , the vector subspace M of E generated by \(H \cap H'\) and a is a hyperplane; since \(H \cap H' \subset L \cap L_1\) and \(a \in L \cap L_1\) , we have \(M \subset L \cap L_1\) , hence \(M = L = L_1\) . Q.E.D.
\end{enumerate}

PROPOSITION 10. Let C be a chamber, let H and H' be two walls of C, and let L be a hyperplane meeting \(D_{H}(C) \cap D_{H'}(C)\) . Assume that H is distinct from H' and that one of the following conditions is satisfied:

\begin{enumerate}
\def\labelenumi{\alph{enumi})}
\item
  The hyperplanes H, H' and L are parallel.
\item
  The hyperplanes H and H' are not parallel, and H ∩ H' ⊂ L.
\end{enumerate}

Then, L meets C.

Let \(b\) (resp. \(b'\) ) be a point of the face of C with support H (resp. H'); it is immediate that every point of the segment \([bb']\) distinct from \(b\) and \(b'\) belongs to C.

Introduce an affine function \(f\) that vanishes at every point of \(H\) and is such that \(f(x) > 0\) for \(x\) in \(D_H(C)\) ; similarly, introduce an affine function \(f'\) having an analogous property with respect to \(H'\) . By applying Lemma 2, we can find numbers \(\lambda\) and \(\lambda'\) and an affine function \(g\) having the properties stated in the lemma. We have \((\lambda, \lambda') \neq (0, 0)\) , and for every point \(x\) of \(L \cap D_H(C) \cap D_{H'}(C)\) , we have \(f(x) > 0\) , \(f'(x) > 0\) and \(\lambda.f(x) + \lambda'.f'(x) = 0\) , so \(\lambda\lambda' < 0\) . On the other hand, we have \(g(b) = \lambda'.f'(b)\) and \(g(b') = \lambda.f(b')\) , and since \(f(b') > 0\) , \(f'(b) > 0\) , we have \(g(b)g(b') < 0\) . The points \(b\) and \(b'\) are thus strictly on opposite sides of the hyperplane \(L\) , hence there exists a point \(c\) of \(L\) which belongs to \([bb']\) , and is distinct from \(b\) and \(b'\) , hence that belongs to \(C\) .

